\section{Conclusion}

This paper evaluated the performance of hyperparameter-optimised machine learning and deep learning models for stock price prediction and Shariah-compliant portfolio optimisation within the BIST Participation All Index. Using daily closing prices of 163 BIST XKTUM stocks from 2007 Q1 to 2022 Q4, the study applied a two-stage empirical framework. In the first stage, LSTM, Bi-LSTM, GRU, XGBoost, Random Forest and SVM models were trained and compared using MAE, MSE and RMSE. In the second stage, the predicted returns from the sequence-based models were integrated into a mean-variance optimisation framework supported by L1 regularisation to examine portfolio-level performance.

The findings show that XGBoost achieved the strongest stock price prediction performance, recording the lowest MAE, MSE and RMSE values among the tested models. This suggests that XGBoost was particularly effective in capturing the non-linear structure of the BIST XKTUM dataset. However, the portfolio optimisation results show that the most accurate forecasting model does not automatically generate the best portfolio outcome. While XGBoost performed best in prediction, the sequence-based models provided the basis for portfolio optimisation because of their ability to model temporal dependencies.

In the portfolio optimisation stage, the LSTM-based portfolio achieved the highest ending equity value, indicating the strongest cumulative portfolio growth. The Bi-LSTM model provided the most favourable risk-return balance, achieving the highest annualised Sharpe ratio. The GRU model produced lower ending equity but remained competitive in terms of volatility and risk-adjusted performance, suggesting resilience under changing market conditions. These findings confirm that model evaluation in financial applications should consider not only prediction accuracy, but also portfolio-level outcomes such as ending equity, volatility and Sharpe ratio.

The paper contributes to Islamic finance and Islamic fintech literature by showing how machine learning and deep learning models can support Shariah-compliant portfolio construction. By applying the models within the BIST XKTUM universe, the study demonstrates that data-driven investment techniques can be used without moving outside an ethically screened equity framework. This is relevant for Islamic robo-advisory, digital wealth management and automated Shariah-compliant investment platforms.

The study also highlights the importance of computational infrastructure in large-scale financial modelling. The use of NVIDIA V100 GPUs enabled efficient hyperparameter tuning and model training across a large stock universe. However, GPU acceleration should be understood as a computational support mechanism rather than a direct explanation for predictive superiority.

Future research may extend this analysis to other Islamic equity indices, compare BIST XKTUM with global Shariah-compliant benchmarks, and incorporate additional data sources such as macroeconomic indicators, firm-level fundamentals, news sentiment and social media sentiment. Further research may also compare L1 regularisation with alternative sparse allocation methods such as Elastic Net, Hierarchical Risk Parity and robust portfolio optimisation. Finally, explainable artificial intelligence methods should be explored to improve transparency, accountability and trust in Islamic fintech applications.